% Business scope (bonus): deployment, integration, continuous training, budget.
\section{Business Scope: Deployment, Integration, Continuous Training and Budget}
\label{sec:business}

\paragraph{What is sold and to whom.} A software analytics layer for a distribution utility (STEG) that sits on
top of the smart-meter infrastructure it already operates: no new device at the customer. It turns the
half-hourly meter data into verdicts that an operator can trust and act on. Its value drivers are fewer wasted
field visits (a genuine grid event is not sent out as a ``faulty meter''), earlier detection of meter
manipulation, correct billing when a meter drifts, and operator time spent only where it matters. These will be
quantified in Phase~2 on the simulated nodes; no saving is claimed here.

\subsection{Deployment model}
The system runs centrally, next to the utility's meter data management system (MDMS): Goals~1--3 run per node
every day (every 30~minutes for urgent rules such as supply loss), Goal~4's coordinator runs once a week. It
needs, per neighbourhood: the existing smart meters, one accurate meter at the head of each low-voltage feeder,
and the operator console (a web application). Roll-out is staged, and each stage must pass its gate before the
next starts:

\begin{table}[H]
\centering\small
\caption{Staged deployment.}
\begin{tabularx}{\textwidth}{@{}>{\raggedright\arraybackslash}p{3.0cm}>{\raggedright\arraybackslash}p{1.6cm}LL@{}}
\toprule\hdr
\textbf{Stage} & \textbf{Duration} & \textbf{What runs} & \textbf{Gate to the next stage} \\
\midrule
A. Shadow pilot, 1 neighbourhood & 3 months & Goals 1--3 produce verdicts and routes, but nothing is acted on; verdicts are compared with field inspections & False positives within the alarm budget; zero unsafe automatic actions; operators judge explanations useful \\
\addlinespace
B. Live pilot, 3 neighbourhoods & 6 months & Operator workflow live (G0--G3); Goal 4 sharing between the three nodes & Missed-escalation rate below target; canary adoption working; privacy review passed \\
\addlinespace
C. Scale-up, district / region & after Stage B, no fixed end & More nodes join; Goal 4 gains participants, which also strengthens robust aggregation & Review load per operator stays within capacity \\
\bottomrule
\end{tabularx}
\end{table}

\subsection{Integration}
\begin{table}[H]
\centering\small
\caption{Interfaces with the utility's systems. Integration is the main one-off cost.}
\begin{tabularx}{\textwidth}{@{}>{\raggedright\arraybackslash}p{3.2cm}LL@{}}
\toprule\hdr
\textbf{Interface} & \textbf{Content} & \textbf{Mechanism} \\
\midrule
Meter data in & 30-minute interval readings and meter event logs & MDMS export (IEC 61968-9 / CIM messages, or scheduled batch files) \\
Feeder-head data in & Feeder energy and voltage & MDMS or SCADA; a feeder-head meter is installed where none exists \\
Context in & Temperature, irradiance, calendar & National weather service feed; holiday and Ramadan calendar \\
Outage management (out) & Genuine supply-loss events & Event message to the OMS \\
Field service (out) & Inspection requests for suspected meters & Work order in the field-service system \\
Billing (out) & Hold or estimate for readings judged untrustworthy & Flag in the billing workflow (only after G2b/G3 decisions) \\
Identity and audit & Operator roles, two-person approval, hash-chained log & Utility single sign-on; log kept in the utility's archive \\
\bottomrule
\end{tabularx}
\end{table}

\subsection{Continuous training}
\begin{itemize}
  \item \textbf{Daily, automatic:} each meter's 28-day profile and each feeder's balance reference move forward.
  \item \textbf{Monthly:} thresholds and the confidence calibration are refitted on newly confirmed cases (field
        inspection results, operator closures, blind operator judgements), with the event-grouped validation of
        Section~\ref{sec:data}. A new version is promoted only after a shadow comparison with the current one, and
        every version (\code{configs/thresholds.yaml}) can be rolled back.
  \item \textbf{Quarterly:} the learned detectors (Isolation Forest, One-Class SVM, autoencoder) are retrained on
        the latest clean data.
  \item \textbf{Weekly:} Goal~4 sharing rounds; each shared artifact passes the local canary test and operator
        approval before adoption.
  \item \textbf{Monitoring, continuously:} false positives against the alarm budget, calibration (Brier score per
        class), share of reviewed cases per operator-hour, and alignment between AI and operator confidence; a
        breach triggers an early refit or a rollback.
\end{itemize}

\subsection{Budget}
\paragraph{What we measured.} Compute is not a cost driver. The complete Goal~1 experiment on one neighbourhood
for a whole year (30 meters, 525\,600 readings, five-fold training of six detectors) runs in about 30 seconds on a
laptop. Daily operation scores only one day per node, so a single ordinary server can serve many neighbourhoods;
explanations are generated only for flagged meter-days. The real operating cost is \textbf{operator review time}:
at the current false-positive rate (2.3\% of clean meter-days, Section~\ref{g1-sec:goal1}), a 30-meter
neighbourhood produces about 250 false-alarm meter-days a year, most of which the routing keeps at G0--G1; reducing
this is a Phase~2 objective.

\paragraph{Pilot budget (stages A and B, 9 months).} Effort is given in person-months; the unit costs are
\emph{planning assumptions} to be validated with STEG, not quotes, and the total scales linearly with them.

\begin{table}[H]
\centering\small
\caption{Indicative pilot budget. Unit costs are planning assumptions (TND = Tunisian dinar).}
\begin{tabularx}{\textwidth}{@{}L>{\raggedleft\arraybackslash}p{2.6cm}>{\raggedleft\arraybackslash}p{3.0cm}>{\raggedleft\arraybackslash}p{2.4cm}@{}}
\toprule\hdr
\textbf{Item} & \textbf{Quantity} & \textbf{Assumed unit cost} & \textbf{Total (TND)} \\
\midrule
Data / ML engineers (Goals 1--4, retraining) & 2 $\times$ 9 person-months & 5\,000 / person-month & 90\,000 \\
Integration engineer (MDMS, OMS, field service, billing) & 6 person-months & 5\,000 / person-month & 30\,000 \\
Distribution-grid expert (rules, scenarios, operator training) & 3 person-months & 5\,000 / person-month & 15\,000 \\
Feeder-head meters where missing (3 nodes, 7 feeders) & up to 7 & 2\,500 installed & up to 17\,500 \\
Server (central, one virtual machine) & 9 months & 500 / month & 4\,500 \\
Operator time for pilot reviews and training & 2 operators, part time & internal & internal \\
\midrule
\textbf{Total (excluding internal operator time)} & & & \textbf{about 157\,000} \\
\bottomrule
\end{tabularx}
\end{table}

\paragraph{Scale-up costs.} Integration is paid once; each new neighbourhood then adds mainly feeder-head meters
where missing, a short shadow period, and operator review time. The language model for explanations can run
locally (no per-call fee) or as a service; in both cases only flagged meter-days are explained, and the plain
evidence summary is always available as a free fallback.
